\newif\ifglobal

%%%%%%%%%%%%%%%%%%%%%%%
%%% Compile to view %%%
%%%%%%%%%%%%%%%%%%%%%%%

%%%%%%%%%%%%%%%%%%%%%%%%%%%%%%%%%%%%%%%%%%%%%%%%%%%
%%% Readme:                                     %%%
%%% https://www.overleaf.com/read/bwdjvfjksvjy/ %%%
%%%%%%%%%%%%%%%%%%%%%%%%%%%%%%%%%%%%%%%%%%%%%%%%%%%

% >>> M?: marks a module setting number or important notice

% >>> M4: Set {./relative-path} to external files for this module <<<
\def \modulefiles {./64-ANACOM}
% To add an external file, use:
% (Replace '---' with actual filename)
% '\input{\modulefiles/tables/---.tex}' for tables
% '\includegraphics{\modulefiles/figures/---}' for figures
% '\subfile{\modulefiles/---}' for register files

% >>> M1: This module is in global project? <<<
% 'YES' - uncomment; 'NO' - comment out
\globaltrue

\ifglobal
    \documentclass[main\_\_CN.tex]{subfiles}
   
\else
    % Font "FandolSong-Regular" used by ctex does not have GSUB table.
    % It causes warnings that do not affect compilation and can be ignored.
    \PassOptionsToPackage{quiet}{fontspec}

    \documentclass[openany, 10pt]{book}
   
    % >>> M2: In ./chip-spec-settings/preamble-trm-module, also find
    %         settings for visibility of todo notes, watermark <<< 
    \usepackage{preamble-trm-module}


    % Enable Chinese typesetting
    \usepackage{ctex}

    % >>> M3: Temporary labels in Chinese
    %         you can add your own labels there <<<
    \myexternaldocument{temp-labels-cn}

\fi %\ifglobal

\setcounter{secnumdepth}{4}

% >>> M5: If this module needs any updates in
% './chip-spec-settings/preamble-trm-module.sty'
% (include additional package, etc.), contact your module PM
% or create the file './chip-spec-settings/changelog.tex' make the
% necessary updates yourself and document them in this file <<<

\raggedright


\begin{document}


% Sets language into which pre-defined bits of text are translated
\selectlanguage{chinese}

% Do not delete this block
\ifglobal
% Add the following front matter if
% this module is outside of global project
\else
    \InputIfFileExists{readme}{}{}
    {\let\clearpage\relax \listoftodos}
    {\let\clearpage\relax \listofdonetodos}
    \tableofcontents
    %\listoftables
    %\listoffigures
\fi


%%%%%%%%%%%%%%%%%%%%%%%%%
%%% TRM Module Begins %%%
%%%%%%%%%%%%%%%%%%%%%%%%%

\hypertarget{anacom}{}
\chapter{模拟电压比较器}
\label{mod:anacom}

%%%%%%%%%%%%%%%%%%%%
%%% Introduction %%%
%%%%%%%%%%%%%%%%%%%%

\section{概述}

\chipname{} 芯片集成了两个模拟电压比较器，即模拟电压比较器 \regindex{0} 和模拟电压比较器 \regindex{1}。模拟电压比较器依靠支持电压比较功能的特殊芯片焊盘 (PAD) 实现，用于监测 PAD 上的电压变化。每个模拟电压比较器都有外部两个电压接入点，分为主电压和参考电压。模拟电压比较器输出的电压比较结果可以作为 ETM（事件任务矩阵）事件来驱动其他外设的 ETM 任务，也可以触发中断。

\begin{tiplisting}
本章节提供的信息适用于模拟电压比较器 \regindex{0} 和模拟电压比较器 \regindex{1}。除非另有说明，本章节中的模拟电压比较器同时指代模拟电压比较器 \regindex{0} 和模拟电压比较器 \regindex{1}。
\end{tiplisting}
\vspace{-2em}

%%%%%%%%%%%%%%%%%%%%
%%% Feature List %%%
%%%%%%%%%%%%%%%%%%%%

\section{主要特性}
\label{sec:anacom-features}

模拟电压比较器具有以下主要特性：

\begin{itemize}
    \item 支持电压比较功能
    \begin{itemize}
        \item 电压比较模式可配置
        \item 内部参考电压值可配置
    \end{itemize}
    \item 支持电压比较中断
    \item 支持 ETM 事件
\end{itemize}


%%%%%%%%%%%%%%%%%%%%%%%%%%%%%%
%%% Architectural Overview %%%
%%%%%%%%%%%%%%%%%%%%%%%%%%%%%%

\section{系统架构}
\label{sec:anacom-arch-overview}

图 \ref{fig:anacom-architecture} 展示了模拟电压比较器的系统架构。

\begin{figure}[H]
    \centering
    \includegraphics[width=0.6\textwidth]{\modulefiles/figures/anacom_arch_20231018__CN}
    \caption{模拟电压比较器系统框图}
    \label{fig:anacom-architecture}
\end{figure}

如上图所示，模拟电压比较器的结构如下：

\begin{itemize}
    \item 两个 PAD 接口：
        \begin{itemize}
            \item PAD1：主电压接口
            \item PAD0：外部参考电压接口
        \end{itemize}
    \item VREF：内部参考电压接口
    \item mode\_control：寄存器控制的选择开关，用于选择参考电压的输入源
    \item COMP\_OUT：比较器输出的电压比较结果，为片上内部信号，外部无法直接获取
    \begin{itemize}
            \item 若主电压高于参考电压，COMP\_OUT 信号输出为高电平
            \item 若主电压低于参考电压，COMP\_OUT 信号输出为低电平
        \end{itemize}
\end{itemize}


%%%%%%%%%%%%%%%%%%%%%%%%%%%%%%
%%% Functional Description %%%
%%%%%%%%%%%%%%%%%%%%%%%%%%%%%%

\section{功能描述}
\label{sec:anacom-func-descr}

模拟电压比较器具有以下功能：

\begin{itemize}

    \item {\bfseries 电压比较}
    
          模拟电压比较器采用的是支持电压比较功能的特殊 PAD，这种特殊 PAD 由两个通用 PAD 组成。下表列出了模拟电压比较器的 PAD 与 GPIO 管脚之间的映射关系。
          
          \begin{longtable}{ | L{5cm} | L{2cm} | L{2cm} | }
          \caption{模拟电压比较器 PAD 和 GPIO 之间的映射关系} \label{tab:anacom-pad-gpio-mapping}\\ \hline
          \rowcolor{lightgray}
          \textbf{模拟电压比较器} & \textbf{PAD0}  &  \textbf{PAD1}\\ \hline
          模拟电压比较器 \regindex{0} & GPIO51  &  GPIO52  \\ \hline
          模拟电压比较器 \regindex{1} & GPIO53  & GPIO54   \\ \hline
          \end{longtable}
        
          电压比较 PAD 的比较模式和参考电压可配置，具体配置流程见小节 \ref{sec:anacom-prog-procedures} \nameref{sec:anacom-prog-procedures}。
        
          电压比较结果输出为 COMP\_OUT 信号。
    
    \item {\bfseries 电压比较模式控制}

          模拟电压比较器将主电压与参考电压进行比较，参考电压可选内部参考电压和外部参考电压。根据参考电压的选择不同，电压比较模式分为两种：

          \begin{itemize}
            \item 主电压与外部参考电压比较；
            \item 主电压与内部参考电压比较；
          \end{itemize}
        
          电压比较模式可通过 \hyperref[fielddesc:LPSYSREGMODECOMP0]{LPSYSREG\_MODE\_COMP\regindex{n}} (\regindex{n} = 0 \verb+~+ 1) 配置。
    
    \item {\bfseries 参考电压配置}

          内部参考电压范围为 0 \verb+~+ (0.7 × VDD\_IO\_6) V，步长为 0.1 × VDD\_IO\_6 V，其中 VDD\_IO\_6 为模拟电压比较器的供电电压，等于芯片的电源电压（通常为 3.3 V）。内部参考电压值可通过寄存器 \hyperref[fielddesc:LPSYSREGDREFCOMP0]{LPSYSREG\_DREF\_COMP\regindex{n}} (\regindex{n} = 0 \verb+~+ 1) 配置。
        
          外部参考电压直接从 PAD 接入，无需配置，输入范围与内部参考电压一致，为 0 \verb+~+ (0.7 x VDD\_IO\_6) V。

    \item {\bfseries 电压比较中断处理}
        
          电压比较的输出值 COMP\_OUT 信号为高电平或低电平，只要输出值发生变化，就会产生对应的中断信号。

\end{itemize}


%%%%%%%%%%%%%%%%%%%%%%%%%%%%%%%%%
%%% Event Task Matrix Feature %%%
%%%%%%%%%%%%%%%%%%%%%%%%%%%%%%%%%

\section{事件任务矩阵}
\label{sec:anacom-etm-func}

在 \chipname{} 中，模拟电压比较器支持 ETM 功能，即可以通过模拟电压比较器的 ETM 事件触发任意外设的 ETM 任务。关于 ETM 更多详细信息，请参考章节 \ref{mod:evntaskmatrix} \nameref{mod:evntaskmatrix}。这里仅介绍与模拟电压比较器相关的 ETM 事件。

模拟电压比较器可产生的 ETM 事件有：

\begin{itemize}
    \item GPIO\_EVT\_ZERO\_DET\_POS\regindex{n}：表示 COMP\_OUT 信号由低电平变为高电平，即主电压从低于参考电压变为高于参考电压。
    \item GPIO\_EVT\_ZERO\_DET\_NEG\regindex{n}：表示 COMP\_OUT 信号由高电平变为低电平，即主电压从高于参考电压变为低于参考电压。
\end{itemize}

模拟电压比较器不支持任何 ETM 任务。


%%%%%%%%%%%%%%%%%%
%%% Interrupts %%%
%%%%%%%%%%%%%%%%%%

\section{中断}

\chipname{} 的模拟电压比较器可以生成 GPIO\_PAD\_COMP\_INT 中断信号并发送给\hyperref[mod:intmtrx]{中断矩阵}。

这一中断信号由模拟电压比较器的内部中断源生成。表 \ref{tab:anacom-int-sources} 列出了模拟电压比较器的内部中断源、中断触发条件、以及生成的中断信号。

\begin{longtable}{ | L{5.5cm} | L{6cm} | L{4cm} | }
\caption{模拟电压比较器的内部中断源} \label{tab:anacom-int-sources}
\\ \hline
\rowcolor{lightgray}
\textbf{内部中断源} & \textbf{触发条件}  &  \textbf{中断信号}\\ \hline
\label{int:GPIOCOMPnALLINT}GPIO\_COMP\regindex{n}\_ALL\_INT (\regindex{n} = 0 \verb+~+ 1) & COMP\_OUT 信号发生跳变  &  GPIO\_PAD\_COMP\_INT  \\ \hline
\label{int:GPIOCOMPnNEGINT}GPIO\_COMP\regindex{n}\_NEG\_INT (\regindex{n} = 0 \verb+~+ 1) & COMP\_OUT 信号由高电平变为低电平   & GPIO\_PAD\_COMP\_INT \\ \hline
\label{int:GPIOCOMPnPOSINT}GPIO\_COMP\regindex{n}\_POS\_INT (\regindex{n} = 0 \verb+~+ 1)& COMP\_OUT 信号由低电平变为高电平  & GPIO\_PAD\_COMP\_INT \\ \hline
\end{longtable}

\begin{tiplisting}
中断、中断源、中断信号等术语的解释以及它们之间的联系请参考章节 \ref{mod:intmtrx} \textit{\nameref{mod:intmtrx}} > \ref{sec:intmtx-int-terminology} \textit{\nameref{sec:intmtx-int-terminology}}。
\end{tiplisting}

每个中断源都有一组通用的配置寄存器，详见章节 \hyperref[interrupt-config-registers]{\textit{中断配置寄存器}}。具体寄存器名称请查看 \ref{sec:hp-gpio-reg-summ} \textit{\nameref{sec:hp-gpio-reg-summ}}。


%%%%%%%%%%%%%%%%%%%%%%%%%%%%%%
%%% Programming Procedures %%%
%%%%%%%%%%%%%%%%%%%%%%%%%%%%%%

\section{配置流程}
\label{sec:anacom-prog-procedures}

软件配置电压比较器功能的流程如下：

\begin{tiplisting}
有关本节提及的寄存器字段的更多信息，请参阅章节 \ref{mod:sysreg} \nameref{mod:sysreg}。
\end{tiplisting}

\begin{enumerate}
    \item 将 \hyperref[fielddesc:LPSYSREGXPDCOMP0]{LPSYSREG\_XPD\_COMP\regindex{n}} (\regindex{n} = 0 \verb+~+ 1) 设置为 1，启用电压比较器功能。
    \item 禁用 PAD1 的正常数字 PAD 功能，包括 IE、OE、WPU 和 WPD，如果使用外部参考电压，则同时仅用 PAD2 的正常数字 PAD 功能。具体的配置方式请参考章节 \ref{mod:iomuxgpio} \nameref{mod:iomuxgpio}。
    \item 配置比较模式，将 \hyperref[fielddesc:LPSYSREGMODECOMP0]{LPSYSREG\_MODE\_COMP\regindex{n}} (\regindex{n} = 0 \verb+~+ 1) 设置为：
        \begin{itemize}
            \item 0：选择主电压和内部参考电压进行比较；
            \item 1：选择主电压和外部参考电压进行比较。
        \end{itemize}
    \item 通过 \hyperref[fielddesc:LPSYSREGDREFCOMP0]{LPSYSREG\_DREF\_COMP\regindex{n}} (\regindex{n} = 0 \verb+~+ 1) 配置内部参考电压，可调电压范围为 0 \verb+~+ (0.7 × VDD\_IO\_6) V，步长为 0.1 × VDD\_IO\_6 V。
\end{enumerate}

\end{document}
